\ours{} rests on one object: the interpretation rule
$B$ that connects visual observation to task judgment.
Competing explanation chains make these rules explicit
and test them; persistent state organizes their revision
evidence, current validity, and applicability scope, so
that they participate in later interaction.
This section develops the two mechanisms in turn after
fixing the task setting.

\subsection{Task Modeling and Overall Framework}
\label{sec:method:setting}

Consider an agent executing a user task $q$. At step $t$
the agent receives an environment observation $o_t$ and,
given the interaction history $H_t$, proposes an action
$\widetilde{a}_t$. When the proposal depends on chart
understanding, let $I_t$ denote the complete chart
observation relevant to the current judgment; if the
chart has left the current page, $I_t$ refers to the
chart observation actually obtained earlier in the task
that still bears on the current judgment. We express a
chart-driven interpretation as
\begin{equation}\label{eq:o-b-c}
O \land B \Rightarrow C_q(a),
\end{equation}
where $O$ is an observation in the chart, $B$ is the
interpretation rule required to convert the observation
into task semantics, and $C_q(a)$ is the proposition
that the action, or high-level choice, $a$ meets the
requirements of task $q$.
The representation separates a visible phenomenon from
the condition under which it is interpretable: the
observation that a curve moves upward within an interval,
jointly with the rule that in the current axis a higher
position means a larger value, supports the conclusion
that the interval grows and that the corresponding action
is admissible.
For a geometric distortion, the observation may be that
object A's bar stands higher than object B's, while the
interpretation rule is that the current bar-height
ordering preserves the value ordering; if the verified
numeric relations in the figure disagree with the
geometric ordering, it is precisely the applicability of
this rule that is challenged first.
Both $O$ and $B$ may carry multiple items; they are
expressed in structured natural language covering entity
binding, numeric relations, encoding direction, time
range, and comparison conditions.
Relations that can be computed explicitly, such as
numeric comparisons and orderings, are checked with
program assistance; the remaining semantic judgments are
verified jointly with the public task requirements.

The method maintains a within-task rule state $M_t^B$
and, for chart-dependent judgments, runs the following
process:
\begin{equation}\label{eq:process}
\text{current interpretation}
\rightarrow
\text{competing-explanation generation}
\rightarrow
\text{in-chart applicability verification}
\rightarrow
\text{rule-state update}
\rightarrow
\text{action re-derivation}.
\end{equation}
Later steps read the applicable rules from the current
context and restart competing verification only when a
relevant judgment is first formed, when new evidence
conflicts with an existing rule, or when the
applicability conditions of a rule are in question.
The process involves two intertwined time scales.
Local verification decides whether a particular
interpretation rule deserves acceptance or revision;
within-task maintenance decides in which later steps
that rule remains effective.
The parameters of the underlying model stay frozen
throughout, and the rule state evolves with the
observations and verification results actually obtained.

\subsection{Competing Explanation Chains}
\label{sec:method:chains}

\subsubsection{Making the Current Explanation Explicit}
When the agent proposes a chart-dependent choice
$\widetilde{a}_t$, the interpretation module forms, from
the complete chart, the public task, the admissible
history, and the current rule state,
\begin{equation}\label{eq:pi0}
\pi_0 = \langle O_0, B_0, p_0\rangle,
\qquad
p_0 = C_q(\widetilde{a}_t),
\end{equation}
stating separately which phenomena the chart is claimed
to contain, which rules carry the observation to the
conclusion, and which evidence those rules already hold
or lack.
Where the current choice lacks support, the module may
record an interpretation gap directly.
A statement such as ``the color is darker'' can stand as
an observation, whereas ``the risk is higher'' stands
only if a mapping from color to risk holds; making the
chain explicit turns that omitted inference condition
into an object for later verification.
The chain is a checkable record of the task argument:
it exposes which conditions the current choice requires
and directs the counter-questions that follow.

\subsubsection{Counter-Questions and Competing
Explanations}
The competing module poses three connected questions
around the current chain: whether the current
observation suffices for the conclusion, which
interpretation rules the inference depends on, and
whether the same chart admits another evidence-based
explanation under which the current conclusion fails,
or under which the current action continues to hold on
different grounds.
From these questions the module generates a bounded
candidate set
\begin{equation}\label{eq:pis}
\Pi_t = \{\pi_0, \pi_1, \ldots, \pi_K\},
\qquad
\pi_i = \langle O_i, B_i, p_i\rangle,
\end{equation}
where each conclusion $p_i$ may support the current
action, refute it, or support a different task judgment.
All candidates refer to the same complete chart, and the
counter-questions act on the interpretive conditions
between observation and conclusion.
A competing chain is meaningful only if it differs
substantially in observation correspondence,
interpretation rule, applicability condition, or scope
of derivation: two chains that both accept ``the curve
rises'' but adopt forward and reversed value mappings
disagree precisely on the applicability of $B$.
A reversed interpretation enters first as a conditional
hypothesis---if the position--value relation of the
current axis differs from the customary direction, the
original trend judgment must change---and remains
unverified until the chart's scales and correspondences
support the condition.
The competing module may return insufficient support and
is not required to produce an opposing conclusion for
every task.

\subsubsection{Normalization and Shared Rules}
Generated chains are normalized over observation, rule,
applicability scope, conclusion, and evidence
dependence.
Explanations that differ only in wording are merged, and
a rule used in several chains shares one verification
state; different entities or comparison relations keep
their separate evidence.
``Darker means higher risk'' and ``a deeper color stands
for a more severe risk'' map to one rule when object and
scope coincide, whereas comparing Illinois with Kansas
and Illinois with Delaware under the same legend are two
comparisons that shared rules do not merge.
Whereas self-consistency aggregates a distribution of
answers over repeated reasoning paths
\citep{wang2022selfconsistency}, the object of
verification here is the rule each path depends on and
the evidence for it: repeated generations of the same
rule add no evidential weight, and a new explanation
changes later verification only by contributing
different evidence or exposing a new applicability
problem.

\subsection{In-Chart Applicability Verification}
\label{sec:method:verify}

The verification module first extracts the disputes that
genuinely affect conclusions among the candidate chains.
If several chains agree on the observation and disagree
only on mapping direction, verification concentrates on
that mapping; if the dispute concerns entity--value
correspondence, observation binding is checked before
rule applicability is judged.
For each chain the module outputs
\begin{equation}\label{eq:vi}
V_i =
\bigl(\,
v_i^{O},\,
v_i^{B},\,
v_i^{\Rightarrow},\,
E_i
\bigr),
\end{equation}
where each verification status takes the values
support, refute, or undetermined, and $E_i$ records the
in-chart evidence.
Observation $O_i$ is checked for presence in the chart
and for the correctness of entity, value, time, and unit
binding; rule $B_i$ is checked for support by the
current chart and satisfaction of its applicability
conditions; the derivation
$O_i\land B_i\Rightarrow p_i$ is checked for
sufficiency of the confirmed premises, including
candidate coverage, range, and task conditions.
Verification always proceeds against the complete chart,
and the evidence record contains the actual observation
identifier, the content read, and which rule the content
supports or refutes.
To keep the judgment focused on evidence, verification
takes the normalized dispute propositions as input, and
the source role or repetition count of a candidate is
never used as a credibility signal; the module confirms
observations and interpretive conditions before
evaluating conclusions, so that a preferred answer is
not used to justify a rule retroactively.

When a chain is challenged, the method distinguishes
three objects of update.
An \emph{observation error}---a misread number, a wrong
entity binding, a wrong range---revises the affected $O$
and re-verifies the interpretations that depend on it;
unreliable observations cannot be charged to $B$
directly.
An \emph{inapplicable rule} is one whose observation
stands while the mapping from observation to semantics
disagrees with the current chart; such a rule is
revoked, or its applicability narrowed, within the
current scope.
\emph{Insufficient derivation} leaves both observation
and rule supported while the full task conclusion
remains unproven; the module then adds the missing
comparison or keeps an undetermined status, as when one
higher candidate refutes the claim that the current
candidate is highest without yet confirming that the
alternative is globally highest.
This separation keeps the defense centered on
interpretation rules while preserving correct handling
of perception, scope, and task-derivation problems.

Conflicting conclusions between candidate chains
establish only that a conflict exists.
A chain is excluded only when chart-supported evidence
negates its observation, its rule applicability, or its
derivation conditions.
For rules that are mutually incompatible over the same
objects and scope, the method checks which in-chart
information can discriminate them: if two explanations
each retain unresolved evidence, the dispute is kept
open; if they apply to different axes, series, or
ranges, the applicability conditions are split rather
than forced to a single choice.
Verification may support a replacement rule, or only
negate the original one: discovering that bar height and
visible value ordering disagree does not entail that
``shorter means larger,'' and the sound revision may
instead be that bar height in the current figure does
not directly support quantity ordering, which must rely
on verified entity--value relations.
Rule revision therefore includes both replacing a
mapping and disqualifying an inference route within the
current chart.

\subsection{Rule Revision and Persistent State Update}
\label{sec:method:state}

Each interpretation rule is maintained as a record
\begin{equation}\label{eq:m}
m_i = \langle
\mathrm{id}_i,\,
B_i,\,
\sigma_i,\,
s_i,\,
E_i,\,
\Delta_i,\,
v_i
\rangle,
\end{equation}
where $\mathrm{id}_i$ identifies the rule for later
reference, $B_i$ is the current interpretation rule,
$\sigma_i$ is its applicability scope over chart
objects, encoding components, and semantic conditions,
$s_i$ is its status among pending, currently active,
disputed, and revoked, $E_i$ is the observed evidence
supporting or refuting the rule, $\Delta_i$ links the
original rule, the revised rule, and the reason for the
change, and $v_i$ is the revision version.
The persistent state stores how the current chart is to
be read and on what evidence; concrete object choices,
task answers, and click actions are produced by the
task reasoning itself and recorded in the execution
history.
A rule's scope is fixed by the conditions under which it
holds: a position--value mapping binds to the numeric
axis and series concerned, a color reading binds to the
current legend and metric, and a rule about overall
trend may further involve temporal order and comparison
range.
Page changes unrelated to the task goal do not
invalidate a rule by themselves.

Verification produces rule-revision proposals
$\Delta_t^B$, which the state manager applies as
$M_{t+1}^B = U(M_t^B, \Delta_t^B)$.
When the original rule is supported and the new proposal
lacks effective evidence, the original rule is kept and
the proposal is retained as a pending record.
When the original rule is refuted and a replacement is
supported, the original loses applicability in the
affected scope and the replacement version is activated.
When the original is refuted without a settled
replacement, the rule is revoked and the current
interpretation gap is recorded.
When the original and a competing rule remain in
effective conflict, the affected rules enter a disputed
status and verification continues within budget, and
when two rules apply to different objects or scopes they
are stored separately rather than overwritten globally.
The ``currently active'' status expresses the method's
judgment under the evidence at hand; new effective
evidence can change it, and can also show that an
earlier verification was itself mistaken.
After revocation, a rule's content, scope, and reason
remain in the revision history; if the same rule
reappears later, it must answer the recorded refutation
or supply evidence that changes the applicability
judgment, since repetition, later position, or a more
confident phrasing forms no new revision ground by
itself.

A revision record answers why the original rule had to
change, what interpretation is adopted now, and under
which conditions the change applies: for a reversed
linear axis, the record replaces ``higher position,
larger value'' because the visible scale increases
downward, states that an upward move of the bound
series now reads as a decrease, and scopes this to the
current chart, axis, and series---supporting trend
judgments over further intervals without storing a
fixed answer of growth or decline.

An invalidated rule does not by itself condemn the
current action, since
$\lnot B_0 \not\Rightarrow \lnot C_q(a_0)$.
After updating $B$, the method re-examines the current
judgments that depended on it: a judgment that lost its
basis enters re-review and may be retained if other
effective interpretations still support it, so the two
admissible outcomes are a rule change with the action
retained (a wrong reason for a right choice) or rule
and action changing together.
Judgments already recorded in the history are marked as
basis-changed for later re-confirmation; the original
interaction history stays intact, and the currently
effective rules are organized by a separate read
interface.

\subsection{Reading and Using Rules over Continued
Execution}
\label{sec:method:read}

Let the current interpretation context $c_t$ collect the
relevant chart objects, encoding components, metrics,
and task scope.
The state reader returns
$\mathcal{R}_t = \operatorname{Read}(M_t^B, c_t)$, of
which the directly usable part is
\begin{equation}\label{eq:read}
\mathcal{R}_t^{\mathrm{active}}
=
\{\, m_i \in M_t^B :
s_i = \mathrm{active}
\land
\operatorname{Match}(\sigma_i, c_t) = \mathrm{true}
\,\}.
\end{equation}
Scope matching yields applicability, non-applicability,
or undetermined; an undetermined match triggers
verification of the corresponding condition rather than
a default assumption that a stored rule covers the new
chart.
The read result also carries the disputes and revocation
records relevant to the current judgment, so the agent
can separate interpretations still in force from ones
that appeared before and lost support; revoked content
never participates in derivation as a current fact.
When a chart has temporarily left the viewport, the
reader keeps supplying applicable rules through the
chart references established within the task, and new
chart objects or encoding changes are bound and verified
afresh.

The base agent produces its next proposal from the
current observation, history, and rule view,
\begin{equation}\label{eq:policy}
(\widetilde{a}_t, \widetilde{\pi}_t)
=
\mathcal{F}_\theta
\bigl(q, o_t, H_t, \mathcal{R}_t\bigr),
\end{equation}
where $\widetilde{\pi}_t$ is a short interpretation
record of the present judgment.
Rules about the same chart can remain valid while the
correct action changes with task conditions, so the
method reuses interpretation evidence across steps while
re-deriving actions from the current observation and the
public task.
If a new proposal explicitly depends on a rule that is
revoked within its scope, the harness returns the
relevant revision evidence and requires re-derivation;
if the proposal brings new evidence sufficient to
challenge an active rule, competing verification is
entered again.
This treatment both restrains old conventions from
regaining force without grounds and preserves a channel
for correcting earlier verification mistakes.
Ordinary navigation, input, and other actions that do
not involve a change of chart interpretation follow the
original execution interface, and the persistent state
acts mainly through later contextual reads, with
complete competing verification executed only when a
trigger condition holds.

Three observable entry points start verification.
The first chart-dependent judgment of a task expands the
current interpretation into
$O \land B \Rightarrow C_q(a)$ and establishes the
task's rule state.
An interpretation conflict re-verifies the disputed
part whenever a later proposal uses a rule inconsistent
with the effective state, or an existing rule is
challenged with in-chart evidence.
An applicability change checks whether existing rules
still hold when the chart, its encoding, the series
binding, or the task scope changes in ways that can
affect interpretation.
The entry points are determined by actual observations,
task context, and action proposals; task correctness
and hidden scoring never participate in triggering.

After a rule update, the agent re-judges the action
under the currently effective interpretation and records
one of three outcomes.
\emph{Keep} means the current choice continues to be
supported by an effective interpretation and the task
conditions.
\emph{Revise} means the current choice is negated by
relevant evidence and an alternative choice gains the
support the task requires.
\emph{Unresolved} means an interpretive dispute that
could still change the action judgment persists, or
necessary evidence is missing.
Showing that some candidate exceeds the current choice
refutes the claim that the current choice is highest,
for example, while confirming the replacement still
requires the comparison that the public task demands;
the method treats such checks as part of the task-semantic
verification rather than imposing one fixed comparison
formula on all tasks.
The unresolved status continues whole-chart verification
within the remaining budget, and when the budget is
exhausted the unfinished reason is recorded explicitly.
Verification suggestions, actual choice revisions, and
final submissions are recorded separately, and only
operations carried out by the executor constitute action
outcomes.

\subsection{Full Algorithm and Integration}
\label{sec:method:algorithm}

Algorithm~\ref{alg:cers}
(Appendix~\ref{app:algorithm}) summarizes one execution
of \ours{} under the candidate budget $K$ and the
review budget $L$ for a single verification node.
The framework enters an existing agent at three points:
the observation entrance establishes chart references,
the pre-execution point performs verification and rule
revision, and the later-reasoning entrance reads the
currently applicable rules; planning, execution, and
review modules share the rule state within one task,
while action capabilities and the environment interface
follow the original system.
Explanation generation and visual verification can be
carried out by different calls of the same frozen model.
Deterministic code handles structured record keeping,
management of normalized merges, version updates, scope
retrieval, and checks of relations that can be computed
explicitly.
The competing-chain count, review rounds, and trigger
settings are fixed before a run, and the extra calls and
context overhead are charged to the method's cost.
The rule state is isolated per task: the runtime uses
only the chart and interaction information the task
permits, while the paired chart, hidden answers, and
scoring feedback stay on the evaluation side.
Charts, rules, and action records are stored separately
so that later analysis can trace how evidence changed
and whether operations followed.


